\begin{algorithm}[H]
\SetAlgoLined
\setstretch{0.92}
\SetAlFnt{\normalsize}
\SetArgSty{textnormal}
\SetKwInput{KwIn}{Input}
\SetKwInput{KwOut}{Output}
\SetKwInput{KwInit}{Initialize}
\SetKwFor{For}{for}{do}{end for}
\SetKwIF{If}{ElseIf}{Else}{if}{then}{else if}{else}{end if}
\SetKw{KwRet}{Return}
\SetKw{KwCont}{continue}
\caption{双阈值离散厚度枚举算法}
\label{alg:q3-enumeration}
\KwIn{DSC扫描速率$\beta_{\rm DSC}$、放热曲线$(T,p(T))$及问题一校准后的传热模型参数}
\KwInit{候选集合$\mathcal N=\{0,1,2,3\}$，时间步$\Delta t=0.25\ \mathrm{s}$，模态数$M=40$}
\KwOut{各$n\in\mathcal N$的$t_{15}(n),t_{10}(n),t_W(n),t_{\rm eff,15}(n),t_{\rm eff,10}(n)$及最优层数$n_{15}^*,n_{10}^*$}

对每个$n\in\mathcal N$计算$d_3(n)=0.3(1+n)\times10^{-3}\ \mathrm{m}$\;
计算服装质量$m_c(n)=A(\rho_1d_1+\rho_2d_2+\rho_3d_3(n))$\;
计算负重限制时间$t_W(n)=20[40-m_c(n)]$\;
\For{each $n\in\mathcal N$}{
  初始化温度场为$37\,^\circ\mathrm{C}$，不可逆相变进度$\xi\leftarrow0$\;
  \While{$T_1(0,t)>10\,^\circ\mathrm{C}$ and $t<3600\ \mathrm{s}$}{
    更新外侧自然对流边界$-k_3T_{3,x}=2.38[T_3(L,t)+40]^{1.25}$\;
    由DSC数据与当前$\xi$计算功能层有效比热$c_{2,\rm eff}$\;
    推进一个时间步的模态解，并更新模态系数、平均温度和相变进度\;
    $t\leftarrow t+\Delta t$\;
  }
  插值得到$T_1(0,t)$首次达到$15\,^\circ\mathrm{C}$与$10\,^\circ\mathrm{C}$的时刻$t_{15}(n),t_{10}(n)$\;
  计算$t_{\rm eff,\theta}(n)=\min\{t_\theta(n),t_W(n)\}$，$\theta\in\{15,10\}$\;
}
对$\theta\in\{15,10\}$选择$n_\theta^*=\arg\max_{n\in\mathcal N}t_{\rm eff,\theta}(n)$\;
\KwRet{$n_{15}^*,n_{10}^*,d_3(n_\theta^*),t_{\rm eff,15}^*,t_{\rm eff,10}^*$及候选方案结果}\;
\end{algorithm}
